% !TeX root = ../main.tex
\section{Introduction}
\label{sec:intro}

Motivational Interviewing \citep[MI;][]{miller1991mi} is a counselling style for helping
people resolve ambivalence about change. It is defined less by what the counsellor says
than by how they say it: open questions, reflective listening, affirmation of the client's
own strengths, and a refusal to argue for change on the client's behalf. Those properties
make MI a demanding target for a dialogue system. The behaviour to be learned unfolds over
many turns, the payoff is delayed, and expert-annotated transcripts are scarce.

A now-standard way around the data problem is to close the loop with language models: a
large model plays the patient, another large model plays an evaluator that fills in
validated MI questionnaires, and a small model is trained against the resulting score
\citep{yosef2024assessing, baruch2025pto}. Within that loop, the natural response to
delayed payoff is \emph{search}: instead of asking the oracle to score a candidate
therapist reply in isolation, simulate the conversation forward a few turns under the
current policy and score where the reply \emph{leads}. This is the idea behind the
\emph{Preference Tree with Look-Ahead} in our ICLR 2025 workshop paper, and it is the
component that paper credited for its best results: models trained with look-ahead depth
$\K{=}5$ scored higher and varied less than those trained at $\K{=}0$, and the paper
concluded that ``deeper look-ahead configurations yield the most stable and high-scoring
results'' \citep{baruch2025pto}.

Look-ahead is also, by a wide margin, the most expensive part of the method. Every
candidate continuation at every branch point requires $\K$ further simulated turns before
the oracle can score it, so the number of patient-simulator calls scales with (branch
points) $\times$ (candidates) $\times$ $\K$ $\times$ (iterations). If the lever works, that
cost is the price of the result. If it does not, it is the single largest waste in the
pipeline. Establishing which of those is true is therefore worth doing carefully, and it is
what this chapter does.

\paragraph{The question.}
\emph{Does simulating the future of a conversation before scoring a therapist turn produce
a better MI therapist than scoring the turn on its own --- and if it sometimes does and
sometimes does not, what determines which?}

\paragraph{Why three generations.}
The MI training loop described above was built three times over three years, each time as an
independent re-implementation rather than a continuation, and each time with a stronger
apparatus (Table~\ref{tab:generations}). Generation~1 is the published ICLR experiment
(Llama-2-7B therapist, GPT-3.5 patient and oracle). Generation~2 rebuilt it with a smaller
therapist, a stronger oracle, harder patients and a JSON-schema rubric, sweeping four
training oracles. Generation~3 rebuilt it again in higher precision, added a second
optimizer, eight evaluation instruments, persona-level pairing, and a held-out grader from
a different model family. All three ran the same look-ahead comparison, $\K \in \{0,5\}$.
That makes look-ahead the one lever measured identically in all three --- a natural
replication series that has never been read as one.

\paragraph{What we find.}
Read as one series, the answer changes sign. Look-ahead clearly helps in generation~1, does
nothing measurable in generation~2, and never leads in generation~3. Reporting that
honestly requires ruling out two mundane explanations, and both are ruled out here:

\begin{itemize}\itemsep2pt
  \item \textbf{It is not the statistics.} The published comparison was unpaired and
  best-versus-best (the best $\K{=}5$ model against the best $\K{=}0$ model), a design that
  selects on the outcome. Re-analysing the same data persona-paired at matched iteration
  counts makes the look-ahead effect \emph{larger} and more significant, not smaller
  (\S\ref{sec:gen1}). The original finding survives a stricter test than it was given.

  \item \textbf{It is not the grader.} Generation~1 was graded by GPT-3.5 through a
  regex-parsed rubric; generations~2 and~3 by \primaryjudge{} through a JSON-schema rubric.
  We re-scored all $1{,}344$ of generation~1's conversations with the generation-3 oracle,
  and the effect survives essentially intact (\S\ref{sec:regrade}). Generations~2 and~3
  already share that oracle --- verified byte-identical prompts, schemas and labels --- so
  after this one re-scoring pass all three generations sit on a single measurement axis.
\end{itemize}

With those two controls in place, the reversal is a statement about the experimental
settings themselves. We then ask what in those settings predicts it, and find one candidate
that is measurable in every generation: how well the \emph{myopic} ($\K{=}0$) baseline
learns. Where the one-step reward already drives learning, look-ahead adds nothing; where
it does not, look-ahead rescues the run. We report that as an interpretation rather than a
finding, because it rests on five training arms and one of them contradicts it
(\S\ref{sec:discussion}).

\paragraph{Contributions.}
\begin{enumerate}\itemsep3pt
  \item \textbf{A three-generation replication of the look-ahead result on one measurement
  axis}, with the grader difference removed by re-scoring rather than assumed away
  (\S\ref{sec:method}, \S\ref{sec:regrade}).

  \item \textbf{A re-analysis of our own published result that strengthens it}, showing
  that the reversal cannot be attributed to the original paper having used a weak test
  (\S\ref{sec:gen1}).

  \item \textbf{A negative result stated at full strength}: across $105$ persona-paired
  contrasts in generation~2 and $8$ matched iterations under two independent graders in
  generation~3, look-ahead never establishes an advantage, and in generation~3 one
  iteration significantly favours its absence (\S\ref{sec:gen2}, \S\ref{sec:gen3}).

  \item \textbf{A mechanism consistent with all three generations.} Look-ahead rescales the
  reward signal without sharpening it: it multiplies the within-group spread by
  ${\sim}1.55\times$ with margin and standard deviation rising by the identical factor,
  creates no additional branch points, and at matched policy adds no reward faithfulness.
  More look-ahead data does not help either --- an arm carrying $1.7\times$ the preference
  pairs scored no better (\S\ref{sec:mechanism}).
\end{enumerate}

\paragraph{A note on reporting a non-replication of one's own work.}
The result in \S\ref{sec:gen3} contradicts the headline of a paper on which the author is
first author. We report it in the strongest form the evidence supports, for two reasons.
The narrow reason is that the two obvious escape hatches --- weak statistics, changed
grader --- were tested and neither holds, so the disagreement is real and informative. The
broader one is that a method whose benefit depends on a boundary condition is more useful
to a practitioner than a method reported as unconditionally good: \S\ref{sec:discussion}
ends with a cheap diagnostic for deciding whether look-ahead is worth paying for in a given
setup, which the original framing could not have produced.
